% Canonical Paper IV section.
A quotient fiber can grow with word length, but cardinality, information, and
running time are different measurements.  This section records the exact
counting statements that survive without a machine model.

\subsection{Words, image growth, and fiber multiplicity}

Let $\Sigma$ be a finite alphabet with $s=|\Sigma|\ge2$.  For each $n$, let
\[
  \rho_n:\Sigma^n\longrightarrow Y_n
\]
be an evaluation or condensation map.  Define
\[
  N_n=|\rho_n(\Sigma^n)|,
  \qquad
  M_n=\max_{y\in Y_n}|\rho_n^{-1}(y)|.
\]

\begin{proposition}[Fiber--image inequality]
\label{prop:fiber-image-inequality}
For every $n$,
\begin{equation}
  M_n\ge \frac{s^n}{N_n}.
  \label{eq:fiber-image-inequality}
\end{equation}
\end{proposition}

\begin{proof}
The $s^n$ words are partitioned among $N_n$ nonempty fibers.  At least one
fiber has size at least the average.
\end{proof}

The inequality does not require a group, metric, or probability model.  It
shows only that small image growth forces some large history fiber.

\subsection{A conditional-information form}

Let $W_n$ be uniformly distributed on $\Sigma^n$ and put
$Y_n=\rho_n(W_n)$.  Shannon entropy is measured in bits~\cite{Shannon1948}.

\begin{proposition}[Evaluation information loss]
\label{prop:evaluation-information-loss}
One has
\begin{equation}
 H(W_n\mid Y_n)
 =n\log_2s-H(Y_n)
 \ge n\log_2s-\log_2N_n,
 \label{eq:conditional-information-loss}
\end{equation}
and
\[
  H(W_n\mid Y_n)\le\log_2M_n.
\]
\end{proposition}

\begin{proof}
Because $Y_n$ is a deterministic function of $W_n$,
$H(W_n,Y_n)=H(W_n)=n\log_2s$, while the chain rule gives
$H(W_n,Y_n)=H(Y_n)+H(W_n\mid Y_n)$.  Also
$H(Y_n)\le\log_2N_n$.  Conditional on $Y_n=y$, the word lies in a set of size
at most $M_n$, so its conditional entropy is at most $\log_2M_n$; averaging
preserves the bound.
\end{proof}

For a nonuniform distribution, the lost information depends on the fiber
weights, not only their cardinalities.  For an individual word, Kolmogorov or
description complexity would introduce another model and is not being used
here.

\subsection{Word metrics and coarse growth}

Let $G$ be a finitely generated group with finite symmetric generating set
$S$.  Write $|g|_S$ for word length and
\[
  \beta_S(n)=|\{g\in G:|g|_S\le n\}|.
\]

\begin{proposition}[Finite generating sets give comparable growth]
\label{prop:word-metric-comparison}
If $S$ and $T$ are finite generating sets, there are constants $A,B\ge1$ such
that
\[
  |g|_T\le A|g|_S,
  \qquad
  |g|_S\le B|g|_T
\]
for all $g\in G$.  Consequently
\[
  \beta_T(n)\le\beta_S(Bn),
  \qquad
  \beta_S(n)\le\beta_T(An).
\]
\end{proposition}

\begin{proof}
Every generator in one finite set is a bounded-length word in the other.  Let
$A$ and $B$ be the corresponding maximum lengths and substitute each letter in
a shortest word.
\end{proof}

Polynomial versus exponential group growth is therefore a coarse property of
the finitely generated group, not of one particular generating set.  Deep
structure theorems classify polynomial growth~\cite{Gromov1981}; the present
paper needs only the elementary comparison above.

\subsection{Four growth functions that must not be merged}

For a history language and evaluation map one may encounter:
\begin{align*}
  W(n)&=\#\{\text{literal words of length at most }n\},\\
  N_\rho(n)&=\#\{\text{distinct evaluated operators reached by those words}\},\\
  M_\rho(n)&=\max_g\#\{\text{words of length at most }n\text{ evaluating to }g\},\\
  R_C(n)&=\#\{\text{contextual residual classes exposed at a selected cut}\}.
\end{align*}
A search algorithm adds still another quantity: the number of states it
actually explores under a specified strategy.  None of these functions is the
running time by definition.

\begin{example}[Free reduction]
The ball of radius $n$ in a free group of rank at least two grows
exponentially, but a given word can be freely reduced by a stack in time linear
in its input length.  Exponential group growth therefore does not imply
exponential evaluation time.
\end{example}

\begin{example}[Finite projective image]
For fixed $q$, every sufficiently long bilateral word over a finite generating
alphabet evaluates into the finite group $\PGL_2(\F_q)$.  Literal word counts
continue to grow, while $N_\rho(n)$ eventually saturates at at most
$q(q^2-1)$.  The growing fibers measure many realizations of a finite semantic
state, not unbounded state-space size.
\end{example}

\subsection{Geometry requires a comparison map}

Hyperbolic spaces and many groups have exponential ball growth.  A history
language may also have exponential word growth.  Equality of the adjectives
``exponential'' does not provide a quasi-isometry, coding theorem, or
simulation.  A geometric complexity theorem would need at least:

\begin{enumerate}
  \item a map from histories or live states into the geometric space;
  \item upper and lower distortion estimates for the selected metrics;
  \item a simulation relating geometric steps to machine actions;
  \item a representation theorem linking geometric volume to encoded states.
\end{enumerate}

Without these data, curvature and volume remain motivating analogies.  The
fixed-model identities in the case studies below are not promoted to a general
curvature--complexity law.
